% Generated draft: requires mathematical and source review.
\documentclass[11pt]{article}
\usepackage[margin=25mm]{geometry}
\usepackage{fontspec}
\setmainfont{Latin Modern Roman}
\usepackage{amsmath,amssymb,mathtools,unicode-math}
\setmathfont{Latin Modern Math}
\usepackage{enumitem,needspace,xurl,hyperref}
\setlength{\parindent}{0pt}
\setlength{\parskip}{5pt}
\setlength{\emergencystretch}{3em}
\newcommand{\N}{\mathbb{N}}
\newcommand{\Z}{\mathbb{Z}}
\newcommand{\Q}{\mathbb{Q}}
\newcommand{\R}{\mathbb{R}}
\newcommand{\C}{\mathbb{C}}
\newcommand{\F}{\mathbb{F}}
\newcommand{\A}{\mathbb{A}}
\newcommand{\PP}{\mathbb{P}}
\newcommand{\E}{\mathbb{E}}
\newcommand{\eps}{\varepsilon}
\newcommand{\dd}{\,\mathrm{d}}
\begin{document}
\section*{UnsolvedMath ID: 30005994}
\textbf{Source problem code:} catalogue entry
\subsection{Definitions and mathematical statement}
Let \(N\in\Z\) satisfy \(2\leq N\leq6\). Interpret \(B^N\) as the open unit ball \(\{x\in\R^N:|x|<1\}\), and interpret the boundary map \(\operatorname{Id}\) as \(x\mapsto x\). Define the admissible gradient fields by
\[
\mathcal V=\left\{u\in H^1(B^N;\R^N):\ \exists\phi\in H^2(B^N;\R),\ u=\nabla\phi\ \text{a.e. in }B^N,\ \operatorname{Tr}(\nabla\phi)(x)=x\ \text{on }\partial B^N\right\},
\]
where \(\operatorname{Tr}\) is the Sobolev trace and the boundary equality is understood in the trace sense. Uniqueness concerns the vector field \(u\), not its scalar potential, which can be changed by an additive constant.

For a parameter \(\varepsilon\), let the source-designated radial vortex be
\[
u_\varepsilon(x)=f_\varepsilon(|x|)\frac{x}{|x|},\qquad x\in B^N\setminus\{0\}.
\]
The value at the origin does not affect its Sobolev equivalence class; the supplied record does not specify the profile equation, profile regularity, or normalization defining \(f_\varepsilon\).

\textbf{Intended mathematical question.} Is \(u_\varepsilon\) the unique minimizer of the source-designated Ginzburg--Landau energy \(G_\varepsilon\) over \(\mathcal V\)? Explicitly, the intended assertion is that \(u_\varepsilon\in\mathcal V\), that
\[
G_\varepsilon(u_\varepsilon)\leq G_\varepsilon(u)\qquad\text{for every }u\in\mathcal V,
\]
and that
\[
u\in\mathcal V,\quad G_\varepsilon(u)=G_\varepsilon(u_\varepsilon)
\quad\Longrightarrow\quad u=u_\varepsilon\ \text{a.e. in }B^N.
\]
The record does not provide the formula or effective domain of \(G_\varepsilon\), the allowed range or quantification of \(\varepsilon\), or a complete definition of the radial vortex. Thus this expresses the intended minimization and uniqueness question but does not determine a fully specified variational problem. The dimension range above follows the supplied cleaned statement. The accompanying literature assessment instead refers to a question for all \(N\geq2\); that broader range is not substituted here.
\subsection{Short English statement}
For each integer dimension \(2\leq N\leq6\), the question asks whether the designated radial vortex uniquely minimizes the designated Ginzburg--Landau energy among gradients of real \(H^2\) potentials whose gradient equals the identity map on the boundary. The energy, radial profile, and parameter range require clarification. The supplied classification is partially solved; its dated literature assessment has not been independently verified.
\subsection{Sources}
\begin{itemize}
\item Supplied local UnsolvedMath problem registry, permanent ID \(30005994\), title: \emph{Gradient-Constrained Ginzburg--Landau Minimizers}; exact source problem code: \texttt{catalogue entry}. Collection: Oberwolfach Reports Open Problems. The record gives proposed year 2024 and names no individual proposer.
\item Supplied report citation: \emph{Calculus of Variations (2025)}, Oberwolfach Reports, DOI 10.4171/owr/2024/37, pp.\ 2113--2178. \url{https://doi.org/10.4171/owr/2024/37}. This is the authoritative report citation supplied with the record; it has not been independently consulted here.
\item Supplied literature source: Ignat--Nahon--Nguyen, gradient fields. \url{https://arxiv.org/abs/2310.11384}.
\item Supplied literature source: Oberwolfach Report 37/2024. \url{https://ems.press/content/serial-article-files/50045}.
\item Literature assessment dated exactly \texttt{2026-08-22}: the supplied curation labels the problem partially solved and states that uniqueness over gradient fields is proved for \(N\geq4\), while describing the question for all \(N\geq2\) as not fully settled. This is supplied dated curation, not an independently verified literature assessment; the background explicitly requests independent verification of time-sensitive status and cited claims.
\item The registry labels the statement \texttt{corrected\_verified} and reports that it was reconstructed and checked against the official Oberwolfach report. This is a supplied verification claim, not verification performed for this draft.
\end{itemize}
\end{document}
